% © 2026 Ing. David Jaroš — CC BY-NC-ND 4.0
% Canonical covariant-tetrad formulation, authorised 2026-07-15.
% UBT-AI-PROVENANCE-BEGIN
% schema: ubt-ai-provenance/v1
% tier: B_machine_verified
% ai_assistance: disclosed
% human_review: machine-verification
% editorial_responsibility: Ing. David Jaroš
% policy: ../../../../AI_PROVENANCE.md
% notice: Machine-verified against named sources or verifiers; individual attestation is not claimed.
% UBT-AI-PROVENANCE-END

\section{The covariant biquaternionic tetrad $E_\mu$}
\label{sec:canonical:biquaternion_tetrad}

The fundamental UBT field is $\Theta(q,\tau)\in\mathbb B$.  The local geometric
object is its normalized covariant first derivative
\begin{equation}
 \boxed{E_\mu:=\mathcal N_0^{-1/2}D_\mu\Theta,}
 \label{eq:canonical:E-from-Theta}
\end{equation}
where $\mathcal N_0>0$ is a fixed unit-setting constant.  No local denominator
is allowed.

\subsection{Classical Lorentz sector}
Let $i$ be the commuting complex unit and $\mathbf e_k$ the quaternion units.
The real Lorentz slice is
\begin{equation}
 W_L=\{i x^0\mathbf1+x^k\mathbf e_k\mid x^a\in\mathbb R\}.
\end{equation}
Quaternion conjugation is denoted by $\sharp$ and acts as
$(z^0\mathbf1+z^k\mathbf e_k)^\sharp=z^0\mathbf1-z^k\mathbf e_k$.
The classical tetrad sector is defined by $E_\mu\in W_L$ and
$\det(e_\mu{}^a)\neq0$, where
\begin{equation}
 E_\mu=i e_\mu{}^0\mathbf1+e_\mu{}^k\mathbf e_k.
\end{equation}

\subsection{Metric and bivector from one product}
On $W_L$ the symmetrised product is central:
\begin{equation}
 \boxed{\frac12(E_\mu^\sharp E_\nu+E_\nu^\sharp E_\mu)
       =g_{\mu\nu}\mathbf1,}
 \label{eq:canonical:central-metric}
\end{equation}
with
\begin{equation}
 g_{\mu\nu}=e_\mu{}^a e_\nu{}^b\eta_{ab},
 \qquad \eta_{ab}=\operatorname{diag}(-1,1,1,1).
\end{equation}
The metric is therefore not obtained by a trace, real-part projection, phase
projection, or compact-fiber average.  It is the unique coefficient of the
unit element because the full symmetrised product is already central.

The antisymmetric companion
\begin{equation}
 \Sigma_{\mu\nu}:=\frac12(E_\mu^\sharp E_\nu-E_\nu^\sharp E_\mu)
\end{equation}
is a biquaternionic bivector carrying local oriented-plane/spin information.
It is distinct from the connection curvature $[D_\mu,D_\nu]$.

\subsection{Transformation and rank}
Under coordinate changes $E_\mu$ transforms as a covector.  Under local
Lorentz-frame changes, its coefficients transform by a local Lorentz matrix;
Eq.~\eqref{eq:canonical:central-metric} is invariant.

The map $e_\mu{}^a\mapsto g_{\mu\nu}$ has rank ten at every nondegenerate
tetrad.  Its six-dimensional kernel consists of local Lorentz rotations and
boosts.  A proof and exact symbolic verification are given in
\texttt{canonical/gr\_closure/covariant\_tetrad\_rank\_theorem.tex} and
\texttt{tools/verify\_covariant\_tetrad\_rank.py}.
The same $E_\mu$ also has an injective block lift to curved Clifford
operators $\Gamma_\mu$ satisfying
$\{\Gamma_\mu,\Gamma_\nu\}/2=g_{\mu\nu}I_4$; see
\texttt{canonical/geometry/biquaternion\_dirac\_lift.tex}.  No second
spinor-current tetrad is introduced.

\subsection{Status and open bridge}
The local metric rank problem is solved at the tetrad level.  The
metric-compatible frame connection is also kinematically reconstructed from
the tetrad and specified torsion,
$\omega=\mathring\omega(e)+K(T)$; in the torsion-free GR branch this is the
unique Levi--Civita spin connection.  Every constant Lorentz tetrad has the
explicit affine representer
$\Theta=\Theta_0+\sqrt{\mathcal N_0}E_\mu x^\mu$ in the flat inertial branch.

For curved configurations, however, it remains to derive torsion and the exact
paired left/right connection from the UBT action and to solve the implicit
system $E_\mu=D_\mu\Theta/\sqrt{\mathcal N_0}$.  A naive one-sided regular
action is generically flat for invertible $\Theta$ under torsion-free
compatibility; the two-sided bimodule derivative avoids this obstruction but
does not yet prove curved-space existence.
